%%%%%%%%%%  scan idx 455  |  folio 442  %%%%%%%%%%
\textit{Lemme.} Soit \(X'\xrightarrow{f}X\) un morphisme projectif et parfait, avec \(X\) quasi-compact, et soit \(L^\bullet\) un complexe parfait sur \(\cO_{X'}\), acyclique en degrés \(>0\). Alors il existe un entier \(n_0\) tel que pour \(n\ge n_0\) le complexe \(\mathbb Rf_*\bigl(L^\bullet\otimes_{\cO_{X'}}\cO_{X'}(n)\bigr)\) soit acyclique en degrés \(>0\).

On applique ce lemme au complexe réduit à \(\cO_{X'}\); il existe donc un entier \(n_0\) tel que pour \(n\ge n_0\), on ait \(R^if_*(\cO_{X'}(n))=0\) si \(i\ge1\). On redescend alors par récurrence jusqu'à \(n=0\). Pour cela, on considère la suite exacte:
\[
0\longrightarrow
\cO_{X'}(n+1)
\longrightarrow
\cO_{X'}(n)
\longrightarrow
\cO_{Y'}(n)
\longrightarrow0,
\]
et on lui applique le foncteur \(f_*\). Or on a \(\mathbb Rf_*(\cO_{Y'}(n))=i_*\circ\mathbb Rg_*(\cO_{Y'}(n))\). Comme \(Y'=\bfP(N)\), \(\cO_{Y'}(n)\) est \(g_*\)-acyclique pour \(n\ge0\). On a donc pour tout \(n\ge0\), et tout \(i\ge1\):
\[
R^if_*(\cO_{Y'}(n))=0,
\]
et par conséquent les suites exactes:
\[
R^if_*(\cO_{X'}(n+1))
\longrightarrow
R^if_*(\cO_{X'}(n))
\longrightarrow0,
\]
ce qui permet d'achever la récurrence.

ii) Il s'agit de montrer que le morphisme canonique:
\[
\bigoplus_{n\ge0}J^n
\longrightarrow
\bigoplus_{n\ge0}f_*\bigl(\cO_{X'}(n)\bigr)
\]
est un isomorphisme. La propriété étant locale, on peut supposer \(J\) défini par une suite \(f_1,\ldots,f_d\) de sections de \(\cO_X\), telles que l'homomorphisme \((\cO_X)^d\longrightarrow\cO_X\) correspondant soit régulier. Posons: \(S=\bigoplus_{n\ge0}J^n\); on définit un homomorphisme d'algèbres: \(\cO_X[T_1,\ldots,T_d]\xrightarrow{v}S\), surjectif, en envoyant 
